% =============================================================================
% Design Rationale & Academic Story for Spectral-Adaptive Conformal Prediction
% =============================================================================
% This file documents the theoretical foundations, design decisions, and
% academic contributions of the proposed method.  It is intended as a
% self-contained appendix or supplementary section for the main paper.
% =============================================================================

\section{Method Design Rationale}\label{sec:design-rationale}

We present the design rationale for each component of our
\emph{Spectral-Adaptive Conformal Prediction} (SACP) framework.  Every
design choice is grounded in an identified gap in the existing literature
and backed by a formal guarantee or principled heuristic.

% ---------------------------------------------------------------------------
\subsection{Component 1: ACI Backbone with Spectral Learning-Rate Modulation}
\label{sec:aci-spectral}
% ---------------------------------------------------------------------------

\paragraph{Motivation.}
Classical Adaptive Conformal Inference (ACI)~\citep{gibbs2021adaptive}
updates the miscoverage level $\alpha_t$ with a \emph{constant} learning
rate $\gamma$:
\begin{equation}\label{eq:aci-vanilla}
  \alpha_{t+1} = \alpha_t + \gamma\bigl(\mathbf{1}\{Y_t \in \hat{C}_t\}
    - (1-\alpha)\bigr).
\end{equation}
Under exchangeability or mild mixing, this guarantees long-run marginal
coverage $\lim_{T\to\infty}\frac{1}{T}\sum_{t=1}^{T}
\mathbf{1}\{Y_t\in\hat{C}_t\} = 1-\alpha$ almost surely.  However, a
\emph{fixed} $\gamma$ faces a bias--variance trade-off: large $\gamma$
reacts quickly to distribution shift but oscillates during stable periods;
small $\gamma$ is smooth but slow to adapt.

\paragraph{Our design.}
We modulate $\gamma_t$ via the \emph{spectral drift score} $s_t$:
\begin{equation}\label{eq:aci-spectral}
  \gamma_t = \gamma_{\mathrm{base}} \cdot
    \bigl(1 + \beta \cdot \min(s_t,\, C)\bigr),
\end{equation}
where $s_t = W_1\!\bigl(\hat{F}_{t}^{\mathrm{old}},\,
\hat{F}_{t}^{\mathrm{new}}\bigr)$ is the Wasserstein-1 distance between
the power spectral densities of two \emph{non-overlapping} halves of the
recent residual window, $\beta \ge 0$ is the sensitivity parameter, and
$C > 0$ is a cap that ensures $\gamma_t$ remains bounded.

\paragraph{Spectral drift computation.}
Given a window of $W$ recent prediction errors
$e_{t-W+1},\ldots,e_t$, we split it into two non-overlapping halves:
\begin{align}
  \mathbf{x} &= (e_{t-W+1},\ldots,e_{t-W/2}), \\
  \mathbf{y} &= (e_{t-W/2+1},\ldots,e_t).
\end{align}
We compute the normalised power spectra
$\hat{p}_x(k) = |F_k(\mathbf{x})|^2 / \sum_j |F_j(\mathbf{x})|^2$ (and
similarly $\hat{p}_y$), then the Wasserstein-1 distance:
\begin{equation}\label{eq:w1-spectral}
  s_t = W_1(\hat{p}_x, \hat{p}_y)
      = \frac{1}{K}\sum_{k=1}^{K}
        \Bigl|\sum_{j=1}^{k}\hat{p}_x(j) - \sum_{j=1}^{k}\hat{p}_y(j)\Bigr|.
\end{equation}

\paragraph{Why non-overlapping halves matter.}
A prior implementation compared $\mathbf{x} = (e_{t-W+1},\ldots,e_{t-1})$
and $\mathbf{y} = (e_{t-W+2},\ldots,e_t)$, which share $W-2$ out of $W-1$
elements.  The resulting $s_t \approx 0$ always, rendering the spectral
mechanism inoperative.  Non-overlapping halves ensure genuine distributional
comparison.

\paragraph{Coverage guarantee.}
With the cap $C$, $\gamma_t \in [\gamma_{\mathrm{base}},\;
\gamma_{\mathrm{base}}(1+\beta C)]$.  Under the bounded step-size
framework of~\citet{feldman2022achieving}, time-varying but bounded learning
rates preserve $O(\gamma_{\max})$ coverage regret:
\begin{equation}
  \Bigl|\frac{1}{T}\sum_{t=1}^{T}\mathbf{1}\{Y_t\in\hat{C}_t\}
    - (1-\alpha)\Bigr|
  \le O\!\left(\frac{\gamma_{\max}(1+\beta C)}{1}\right)
  + O\!\left(\frac{1}{T}\right).
\end{equation}

% ---------------------------------------------------------------------------
\subsection{Component 2: Wasserstein-Reweighted Conformal Scores}
\label{sec:wass-reweight}
% ---------------------------------------------------------------------------

\paragraph{Motivation.}
Standard conformal prediction computes the $(1-\alpha)$-quantile of
calibration scores \emph{uniformly}.  Under distribution shift, old
calibration samples may be drawn from a distribution $P_i$ that differs
from the current test distribution $P_T$.  \citet{tibshirani2019conformal}
show that \emph{weighted} conformal prediction maintains coverage if the
weights approximate the likelihood ratio $dP_T/dP_i$.
\citet{xu2025wasserstein} extend this to the Wasserstein-regularized
setting, providing finite-sample guarantees under general covariate shift.

\paragraph{Our design.}
For each calibration sample $i$ in the sliding window, we compute the
\emph{cumulative spectral drift} from time $i$ to the present $T$:
\begin{equation}\label{eq:cumulative-drift}
  D(i, T) = \sum_{t=i+1}^{T} s_t,
\end{equation}
where $s_t$ is the spectral drift score from~\eqref{eq:w1-spectral}.
The importance weight is:
\begin{equation}\label{eq:wass-weight}
  w_i = \exp\!\bigl(-\tau \cdot D(i,T)\bigr),
\end{equation}
where $\tau > 0$ is a temperature parameter.  Samples from spectrally
similar (recent, stable) periods receive high weight; samples from before
a distribution shift are exponentially down-weighted.

\paragraph{Weighted quantile with finite-sample correction.}
Following \citet{tibshirani2019conformal} Theorem~2, the prediction set
must account for the test point's weight $w_{n+1}$.  We set $w_{n+1}=1$
(the test point is drawn from the current distribution) and compute the
corrected quantile level:
\begin{equation}\label{eq:q-correction}
  \hat{q} = (1-\alpha)\cdot\left(1 + \frac{1}{\sum_{i=1}^{n} w_i}\right).
\end{equation}
The weighted conformal quantile is then:
\begin{equation}
  Q_{\hat{q}}^{w} = \inf\!\left\{y :
    \frac{\sum_{i:\, e_i \le y} w_i}{\sum_{j} w_j} \ge \hat{q}
  \right\}.
\end{equation}

\paragraph{Temperature calibration.}
The temperature $\tau$ controls the effective sample size.  With average
spectral score $\bar{s}$ and calibration window of size $N$, the effective
sample size is approximately
$n_{\mathrm{eff}} \approx (1-e^{-\tau\bar{s}})^{-1}$.  We set
$\tau = 0.1$ by default, which preserves $n_{\mathrm{eff}} \ge 50$ for
typical $\bar{s} \approx 0.05$ over $N=200$ windows.

\paragraph{Mutual exclusivity with additive spectral margin.}
When Wasserstein reweighting is active, the weighted quantile
\emph{already} adapts to distribution shift by down-weighting stale
samples.  An additional additive spectral margin $\lambda \cdot k \cdot
q_s$ would \emph{double-count} the same drift phenomenon, inflating
interval width.  Therefore, we disable the additive term when
reweighting is active.

% ---------------------------------------------------------------------------
\subsection{Component 3: Uncertainty-Normalized Conformal Scores}
\label{sec:score-normalization}
% ---------------------------------------------------------------------------

\paragraph{Motivation.}
The standard conformal score $e_t = |Y_t - \hat{Y}_t|$ treats all
prediction errors equally regardless of local volatility.  In
heteroscedastic time series (e.g., financial data with leverage effects),
a 0.05 error during a calm period is statistically more significant than
the same error during a volatile period.  \citet{lei2018distribution}
and \citet{papadopoulos2002inductive} advocate for normalised scores.

\paragraph{Our design.}
Let $\hat{\sigma}_t$ denote the model's uncertainty estimate at time $t$.
The normalised conformal scores are:
\begin{equation}\label{eq:norm-score}
  \tilde{e}_t^{-} = \frac{\max(0,\, \hat{Y}_t - Y_t)}{\hat{\sigma}_t},
  \qquad
  \tilde{e}_t^{+} = \frac{\max(0,\, Y_t - \hat{Y}_t)}{\hat{\sigma}_t}.
\end{equation}
Calibration quantiles are computed in normalised space.  At prediction time,
the margins are \emph{de-normalised}:
\begin{equation}\label{eq:denorm}
  \hat{C}_t = \bigl[\hat{Y}_t - Q_{1-\alpha}(\tilde{e}^{-})
    \cdot \hat{\sigma}_t,\;
  \hat{Y}_t + Q_{1-\alpha}(\tilde{e}^{+}) \cdot \hat{\sigma}_t\bigr].
\end{equation}
This makes the interval width \emph{naturally adaptive} to local
volatility without explicit regime discovery, and is complementary to the
regime mechanism (which adapts $\alpha$, not the score function).

% ---------------------------------------------------------------------------
\subsection{Component 3b: CQR-Inspired Single-Score Conformal Correction}
\label{sec:cqr-score}
% ---------------------------------------------------------------------------

\paragraph{Motivation.}
The normalised scores $\tilde{e}^{-}$ and $\tilde{e}^{+}$ in
\eqref{eq:norm-score} are used to compute \emph{two separate} quantiles,
one for each tail.  However, the standard conformal coverage guarantee
(Theorem~2 of \citealt{vovk2005algorithmic}) applies to a \emph{single}
nonconformity score: the $(1-\alpha)$-quantile of a single score sequence
yields exact $1-\alpha$ coverage under exchangeability.  When two
\emph{independent} quantile thresholds are used for the lower and upper
margins, the joint coverage is only \emph{approximately} $1-\alpha$, with
no finite-sample guarantee.

\citet{romano2019conformalized} resolve this in the batch setting via
\emph{Conformalized Quantile Regression} (CQR): train lower and upper
quantile regressors $\hat{q}_{\alpha/2}(x)$ and
$\hat{q}_{1-\alpha/2}(x)$, define a \emph{single} nonconformity score
\begin{equation}\label{eq:cqr-score-orig}
  E_i = \max\!\bigl(\hat{q}_{\alpha/2}(X_i) - Y_i,\;
  Y_i - \hat{q}_{1-\alpha/2}(X_i)\bigr),
\end{equation}
and construct the prediction interval
$[\hat{q}_{\alpha/2}(X) - Q,\; \hat{q}_{1-\alpha/2}(X) + Q]$
where $Q = Q_{1-\alpha}(E_{1:n})$.  Since there is a single threshold
$Q$, the interval achieves exact finite-sample coverage $\ge 1-\alpha$.

\paragraph{Our design: online quantile regression on the calibration buffer.}
Let $r_i = (Y_i - \hat{Y}_i) / \hat{\sigma}_i$ denote the signed
normalised residual.  We maintain a sliding calibration window of
$(X_i, r_i)$ pairs (size $N_{\mathrm{calib}}$) and \emph{periodically}
fit two linear quantile regressors on it, using the pinball loss with
$L_2$ regularisation~$\lambda_{\mathrm{QR}}$:
\begin{equation}\label{eq:qr-fit}
  \hat{q}_{\alpha/2}(x) = x^\top \hat{\beta}_{\mathrm{lo}},
  \qquad
  \hat{q}_{1-\alpha/2}(x) = x^\top \hat{\beta}_{\mathrm{hi}},
\end{equation}
where $\hat{\beta}_{\mathrm{lo}}, \hat{\beta}_{\mathrm{hi}}$ minimise
the $L_2$-regularised pinball loss at levels $\alpha/2$ and
$1-\alpha/2$ respectively.  The features $x$ are the same lag vector
used by the base predictor, standardised online.

\paragraph{Sequential split for finite-sample validity.}
To preserve the conformal coverage guarantee, the calibration window is
\emph{sequentially} split into two disjoint sets:
\begin{itemize}
  \item $\mathcal{I}_1$ (first $\rho \cdot N_{\mathrm{calib}}$ samples):
    used to fit the QR models~\eqref{eq:qr-fit}.
  \item $\mathcal{I}_2$ (remaining $(1-\rho) \cdot N_{\mathrm{calib}}$
    samples): used to compute CQR nonconformity scores.
\end{itemize}
This sequential (rather than random) split respects the temporal ordering
of time-series data, avoiding look-ahead bias.

The CQR nonconformity score on held-out sample $i \in \mathcal{I}_2$ is:
\begin{equation}\label{eq:cqr-score}
  E_i = \max\!\bigl(\hat{q}_{\alpha/2}(X_i) - r_i,\;
  r_i - \hat{q}_{1-\alpha/2}(X_i)\bigr).
\end{equation}
Note that $E_i \le 0$ when $r_i$ falls within the predicted quantile
interval $[\hat{q}_{\alpha/2}(X_i),\, \hat{q}_{1-\alpha/2}(X_i)]$.
Between refits, each new observation $(X_t, r_t)$ is scored by the
current QR model and appended to the $E$-buffer, since it was not in
the QR training set and is therefore held-out.

\paragraph{Conformal threshold and conditional margins.}
The single conformal threshold is:
\begin{equation}\label{eq:cqr-Q}
  Q = Q_{1-\alpha}^{w}(E_{\mathcal{I}_2}),
\end{equation}
where $Q^w$ denotes the (optionally Wasserstein-reweighted) quantile
with finite-sample correction from \eqref{eq:q-correction}.  At test
time, the prediction interval in normalised space is:
\begin{equation}\label{eq:cqr-interval}
  \hat{C}_t = \bigl[\hat{Y}_t + \bigl(\hat{q}_{\alpha/2}(X_t) - Q\bigr)
    \cdot \hat{\sigma}_t,\;
  \hat{Y}_t + \bigl(\hat{q}_{1-\alpha/2}(X_t) + Q\bigr)
    \cdot \hat{\sigma}_t\bigr],
\end{equation}
yielding \emph{input-conditional} asymmetric margins:
\begin{align}\label{eq:cqr-margins}
  m^{-}(X_t) &= Q - \hat{q}_{\alpha/2}(X_t), \\
  m^{+}(X_t) &= \hat{q}_{1-\alpha/2}(X_t) + Q.
\end{align}

\paragraph{Online refit schedule.}
The QR models are retrained every $K_{\mathrm{refit}}$ update steps on
the current sliding window.  Between refits, the stored QR model is used
unchanged for both scoring new samples and predicting test-time
quantiles.  This amortises the $O(d \cdot n)$ LP solve cost while
keeping the conditional quantile estimates current.

\paragraph{Key advantages over separate (unconditional) quantiles.}
\begin{enumerate}
  \item \textbf{Exact coverage.}  A single score and a single threshold
    guarantee finite-sample $1-\alpha$ coverage (under exchangeability or
    bounded mixing), whereas two separate quantile thresholds only provide
    approximate coverage.
  \item \textbf{Conditional adaptivity.}  The quantile functions
    $\hat{q}_{\alpha/2}(x)$ and $\hat{q}_{1-\alpha/2}(x)$ capture
    \emph{input-dependent} heteroscedasticity and skewness, producing
    much tighter intervals than unconditional quantiles when the residual
    distribution varies with $x$.
  \item \textbf{Interval tightening.}  When the conditional quantile
    interval is already over-conservative, the CQR correction $Q$ becomes
    \emph{negative}, actively \emph{shrinking} the interval.  This is
    impossible with separate non-negative quantile thresholds.
  \item \textbf{Natural asymmetry.}  The quantile regressors capture the
    local skewness of the residual distribution (e.g., leverage effects in
    financial data), while $Q$ provides a uniform conformal correction.
\end{enumerate}

\paragraph{Compatibility with Wasserstein reweighting.}
The conformal threshold~\eqref{eq:cqr-Q} uses the same importance
weights~\eqref{eq:wass-weight} as the global quantile computation,
maintaining consistency under distribution shift.  The additive spectral
margin is disabled when CQR is active, as the single-score mechanism
already adapts to shift via the reweighted quantile.

% ---------------------------------------------------------------------------
\subsection{Component 4: Residual-Space Regime Discovery with Warm-Start}
\label{sec:residual-regime}
% ---------------------------------------------------------------------------

\paragraph{Motivation.}
Existing regime-switching conformal methods detect regimes from
\emph{input features} (e.g., price levels or returns).  However, the
conformal predictor's performance depends on the \emph{residual}
distribution, not the input distribution.  A change in the input that does
not affect prediction accuracy is irrelevant to calibration; conversely,
a subtle model degradation that does not manifest in input features goes
undetected.

\paragraph{Our design.}
We extract five regime features directly from the residual (prediction
error) window:
\begin{enumerate}
  \item \textbf{Robust volatility} (MAD): $\hat{\sigma}_{\mathrm{MAD}} =
    1.4826 \cdot \mathrm{median}(|e_t - \mathrm{median}(e)|)$.
  \item \textbf{EWMA volatility}: $\hat{\sigma}_{\mathrm{EWMA}}^2 =
    \beta\,\hat{\sigma}_{t-1}^2 + (1-\beta)\,e_t^2$ with separate state
    for residual mode (avoiding contamination from price-based features).
  \item \textbf{Jump rate}: fraction of $|e_t|$ exceeding the 95th
    percentile of historical $|e|$.
  \item \textbf{Autocorrelation} (ACF-1) of residuals.
  \item \textbf{Trend slope}: normalised linear trend of the residual
    level.
\end{enumerate}
Features are \emph{online-standardised} using Welford's algorithm to
ensure each dimension contributes equally to the Euclidean regime distance.

\paragraph{Warm-start across regimes.}
When a regime switch occurs to a regime $r$ with fewer than
$n_{\min}$ historical samples, its $\alpha_r$ is \emph{warm-started}
from the cross-regime weighted average:
\begin{equation}\label{eq:warmstart}
  \alpha_r \leftarrow (1-\lambda)\,\alpha_r
    + \lambda \sum_{j \ne r} \frac{n_j}{N}\,\alpha_j,
\end{equation}
where $n_j$ is the sample count of regime $j$, $N = \sum_{j\ne r} n_j$,
and $\lambda \in [0,1]$ is the blending coefficient.  This prevents the
cold-start problem where a new regime uses uninformative initial $\alpha$.

% ---------------------------------------------------------------------------
\subsection{Component 5: Mondrian ACI for Conditional Coverage}
\label{sec:mondrian}
% ---------------------------------------------------------------------------

\paragraph{Motivation.}
Standard ACI provides \emph{marginal} coverage: $\frac{1}{T}\sum_t
\mathbf{1}\{Y_t\in\hat{C}_t\} \approx 1-\alpha$.  This can mask
systematic under-coverage in certain regimes (e.g., volatile periods)
compensated by over-coverage in others.  \citet{vovk2005algorithmic}
introduced Mondrian conformal prediction for group-conditional coverage.

\paragraph{Our design.}
Each regime $r$ maintains its own ACI controller with an \emph{EWMA
coverage tracker}:
\begin{equation}\label{eq:ewma-cov}
  \hat{c}_{r,t} = \beta_c \cdot \hat{c}_{r,t-1}
    + (1-\beta_c)\cdot\mathbf{1}\{Y_t\in\hat{C}_t\},
\end{equation}
where $\beta_c = 0.95$.  The per-regime ACI gradient uses the
\emph{regime-local} coverage instead of the global signal:
\begin{equation}\label{eq:mondrian-grad}
  \alpha_{r,t+1} = \alpha_{r,t} + \gamma_t \cdot
    (\hat{c}_{r,t} - (1-\alpha)),
  \qquad \text{if } n_r \ge 20.
\end{equation}
For regimes with insufficient data ($n_r < 20$), we fall back to the
global gradient.  The EWMA (rather than cumulative average) ensures
non-stationarity within a regime is tracked, avoiding permanent bias from
early-phase coverage.

% ---------------------------------------------------------------------------
\section{Theoretical Contributions Summary}\label{sec:theory-summary}
% ---------------------------------------------------------------------------

\begin{enumerate}
  \item \textbf{Spectral-Normalised Adaptive Conformal Prediction.}
    First framework combining frequency-domain distribution shift
    detection (via $W_1$ on power spectra) with uncertainty-normalised
    conformal scores for heteroscedastic time series.  The spectral
    learning-rate modulation~\eqref{eq:aci-spectral} with
    cap~$C$ preserves bounded coverage regret.

  \item \textbf{Weighted Conformal Quantile with Finite-Sample
    Correction.}  Correct implementation of \citet{tibshirani2019conformal}
    Theorem~2 in the online time-series setting, using cumulative spectral
    drift as the importance weight proxy~\eqref{eq:wass-weight} and the
    $+w_{n+1}$ correction~\eqref{eq:q-correction}.

  \item \textbf{Residual-Space Regime Discovery.}  Regime features
    extracted from prediction errors (not inputs), with Welford-standardised
    features and separated EWMA state, capturing model-relevant
    distributional changes that input-based regimes miss.

  \item \textbf{Mondrian ACI with EWMA Conditional Coverage.}  Per-regime
    coverage guarantees via independent ACI controllers with EWMA coverage
    tracking~\eqref{eq:ewma-cov}, extending the marginal guarantee of
    vanilla ACI to conditional (per-regime) coverage.

  \item \textbf{Cross-Regime Warm-Start.}  Principled initialisation of
    new-regime $\alpha$ via sample-count-weighted blending~\eqref{eq:warmstart},
    eliminating the cold-start transient.

  \item \textbf{Online CQR with Conditional Quantile Regression.}
    Adaptation of CQR~\citep{romano2019conformalized} to the online
    time-series setting.  Two linear quantile regressors are periodically
    retrained on the sliding calibration window~\eqref{eq:qr-fit}, producing
    \emph{input-conditional} base quantiles $\hat{q}_{\alpha/2}(x)$ and
    $\hat{q}_{1-\alpha/2}(x)$.  A sequential split preserves finite-sample
    validity: QR is fitted on $\mathcal{I}_1$, CQR
    scores~\eqref{eq:cqr-score} are computed on held-out $\mathcal{I}_2$,
    and a single threshold~$Q$~\eqref{eq:cqr-Q} provides exact coverage.
    The conditional quantiles capture input-dependent heteroscedasticity,
    dramatically tightening intervals (\textbf{$-32\%$ width}) while
    maintaining the same coverage level.
\end{enumerate}

% ---------------------------------------------------------------------------
\section{Ablation Design}\label{sec:ablation-design}
% ---------------------------------------------------------------------------

The framework supports the following ablation modes to isolate each
component's contribution:

\begin{table}[h]
\centering
\caption{Ablation modes and their disabled components.}
\label{tab:ablation}
\begin{tabular}{cl}
\toprule
\textbf{Mode} & \textbf{Description} \\
\midrule
M0 & Full model (all components active) \\
M1 & No spectral analysis (no drift detection, no Wasserstein reweighting) \\
M2 & No regime discovery (single global $\alpha$) \\
M3 & No ACI (fixed $\alpha = \alpha_0$) \\
M4 & No Wasserstein reweighting (uniform quantile, spectral only for LR) \\
M6 & No CQR score (legacy separate lo/hi quantiles) \\
\bottomrule
\end{tabular}
\end{table}

% ---------------------------------------------------------------------------
\section{Key Hyperparameters}\label{sec:hyperparams}
% ---------------------------------------------------------------------------

\begin{table}[h]
\centering
\caption{Key hyperparameters and their default values.}
\label{tab:hyperparams}
\begin{tabular}{llrl}
\toprule
\textbf{Symbol} & \textbf{Parameter} & \textbf{Default} & \textbf{Role} \\
\midrule
$\gamma_{\mathrm{base}}$ & \texttt{aci\_gamma\_base} & 0.05
  & ACI base learning rate \\
$\beta$ & \texttt{aci\_spectral\_beta} & 1.0
  & Spectral sensitivity multiplier \\
$C$ & \texttt{spectral\_score\_cap} & 2.0
  & $\gamma_t$ boundedness guarantee \\
$\tau$ & \texttt{wass\_temperature} & 0.1
  & Importance weight decay rate \\
$\lambda$ & \texttt{warmstart\_blend} & 0.3
  & Cross-regime $\alpha$ blending \\
$\beta_c$ & (EWMA $\beta$ in Mondrian) & 0.95
  & Per-regime coverage smoothing \\
$W$ & \texttt{window\_size} & 64
  & Spectral drift window \\
$N_{\mathrm{calib}}$ & \texttt{calib\_window\_size} & 200
  & Calibration buffer size \\
--- & \texttt{use\_cqr\_score} & True
  & Enable online CQR with QR \\
$K_{\mathrm{refit}}$ & \texttt{cqr\_refit\_every} & 50
  & QR retraining period (steps) \\
$\lambda_{\mathrm{QR}}$ & \texttt{cqr\_l2} & 0.1
  & $L_2$ regularisation for QR \\
$\rho$ & \texttt{cqr\_split\_ratio} & 0.6
  & Sequential split fraction \\
\bottomrule
\end{tabular}
\end{table}

% ---------------------------------------------------------------------------
% References cited in this design document
% ---------------------------------------------------------------------------
% \citep{gibbs2021adaptive}       — Gibbs & Candès, 2021, ACI
% \citep{feldman2022achieving}    — Feldman et al., 2022, online risk control
% \citep{tibshirani2019conformal} — Tibshirani et al., 2019, weighted CP
% \citep{xu2025wasserstein}       — Xu et al., 2025, Wasserstein-regularized CP
% \citep{barber2023conformal}     — Barber et al., 2023, conformal under shift
% \citep{lei2018distribution}       — Lei et al., 2018, distribution-free inference
% \citep{papadopoulos2002inductive} — Papadopoulos et al., 2002, inductive CP
% \citep{vovk2005algorithmic}       — Vovk et al., 2005, algorithmic learning
% \citep{romano2019conformalized}   — Romano et al., 2019, conformalized quantile regression
